%\documentclass{recpad2k}
\documentclass[extendedabs]{recpad2k}

%% Uncomment and set paper number to enable the review ruler
%\recpadreviewcopy{??}

% Caption formatting to match Word template
% 8pt font, color #001A66, regular label/text weight, 6pt space after, centered
\usepackage[font={small,color=captionblue},labelfont=normalfont,textfont=normalfont,skip=6pt,justification=centering,belowskip=0pt]{caption}
% Word template float/caption spacing: figure/table objects use 6 pt
% around the object and captions use 6 pt after the caption paragraph.
\setlength{\floatsep}{6pt plus 1pt minus 1pt}
\setlength{\textfloatsep}{6pt plus 1pt minus 1pt}
\setlength{\intextsep}{6pt plus 1pt minus 1pt}
\setlength{\abovecaptionskip}{6pt}
\setlength{\belowcaptionskip}{0pt}
\definecolor{captionblue}{RGB}{0,26,102}

\title{Author Guidelines for the Portuguese Conference on Pattern Recognition}

% Enter the paper's authors in order
% \addauthor{Name}{email/homepage}{INSTITUTION_CODE}
% INSTITUTION_CODE pode ser um número único (ex: 1) ou uma lista
% separada por vírgulas para múltiplas afiliações (ex: 1,2)
\addauthor{Susan Steel}{ssteel@visioninst.ac.uk}{1}
\addauthor{Petra Moore}{petra.moore@collaborators.com}{2}
\addauthor{Colin Martassen}{colin@collaborators.com}{1,2}
\addauthor{John O'Shea}{ana.ferreira@techlab.pt}{1,3}

% Enter the institutions
% \addinstitution{Name\\Address}
\addinstitution{
 The Vision Institute\\
 University of Borsetshire\\
 Wimbleham, UK
}
\addinstitution{
 Collaborators, Inc.\\
 100 Central Park East,\\
 New York, USA
}
\addinstitution{
 TechLab Research Centre\\
 University of Lisbon\\
 Lisbon, Portugal
}

\runninghead{Steel, Moore, Martassen}{RECPAD Author Guidelines}

% Macro definitions
\def\eg{\emph{e.g}\bmvaOneDot}
\def\Eg{\emph{E.g}\bmvaOneDot}
\def\etal{\emph{et al}\bmvaOneDot}

%-------------------------------------------------------------------------
% Document starts here
\begin{document}

\maketitle

\begin{abstract}
This document outlines the format requirements for papers submitted to RECPAD, the Portuguese Conference on Pattern Recognition. It includes a description of the format and highlights common problems and errors. It also features a greatly simplified way of listing multiple authors and institutions. The abstract should not include citations.
\end{abstract}

%-------------------------------------------------------------------------
\section{Introduction}
\label{sec:intro}
The proceedings of RECPAD will be published in electronic format only. This document illustrates the required document format and includes guidelines for preparing submissions. Papers that do not comply with these requirements may be rejected at any stage of the review process.

\LaTeX\ users should use this template in order to prepare their paper. Users of other packages should emulate the style and layout of this example.  Note that best results will be achieved using {\tt pdflatex}, which is available in most modern distributions.

Although this information is embedded in the document class and style files, the header and headings use Arial (9 pt, except for the title, which is 12 pt), while all other text is in Times New Roman (also 9 pt, except for footnotes and captions, as noted below). The paper title, authors' contacts, headings, captions, and hyperlinks are coloured \#001966 (a dark blue). 

The paper title should be written in Title Case, capitalising the first letter of each word, except for function words (such as \emph{an}, \emph{and}, \emph{as}, \emph{between}, \emph{by}, \emph{for}, \emph{from}, \emph{in}, \emph{of}, \emph{on}, \emph{the}, \emph{through}, \emph{to}, \emph{with}), which are usually not capitalised unless they are the first or last word of the title. Line breaks may be used to properly wrap a title that needs more than one line. Avoid mathematical notation and uncommon special symbols in the title. Standard technical acronyms and widely accepted terminology may be used when necessary.

\subsection{Paper length: two pages, including bibliography and title}
Papers must be no more than 2~pages in length, including the bibliography. Appendices or annexes are not allowed. This strict limit is intended to encourage conciseness.

The bibliography section should begin immediately after the paper text. It may be of any reasonable length. It should not include annotations, figures, or any other additional graphical elements.

\subsection{Citations}

Citations are to be numbered in the order of their appearance, enclosing the citation number in square brackets, for example~\cite{Peterson2002}.

When citing a paper with three or more authors, authors may save space by using ``et~al.'' (not ``et.~al.'' as ``et'' is a complete word). Thus, the following is correct: ``Frobnication has been trendy lately. It was introduced by Peterson~\cite{Peterson2002}, and subsequently developed by Alpher and Fotheringham-Smythe~\cite{AlpherFotheringhamSmythe2003}, and Rudolph
et~al.~\cite{Rudolph2004}.''

This is incorrect: ``\ldots\ subsequently developed by Alpher et~al.~\cite{AlpherFotheringhamSmythe2003}\ldots'' because reference~\cite{AlpherFotheringhamSmythe2003} has just two authors.

%-------------------------------------------------------------------------
\subsection{Footnotes}

Please use footnotes\footnote{This is an example of a footnote. It can span across columns and occupy several lines if necessary.} sparingly. Indeed, try to avoid footnotes altogether and include necessary peripheral observations in the text (possibly within parentheses, as in this sentence). When using footnotes, place them at the bottom of the page where they are referenced. Use 8-pt Times New Roman, single-spaced.

%-------------------------------------------------------------------------
\subsection{Figures and tables}

All figures must be accompanied by a caption placed below the figure, centred, in 8-pt Times New Roman. Captions should be concise and self-contained, allowing the reader to understand the figure without reference to the main text. Label figures sequentially as ``Figure~1'', ``Figure~2'', and so on, followed by a colon and a brief description, ending with a full stop, as demonstrated in Figure~\ref{fig:short}.

\begin{figure}[htpb]
\centering
\fbox{\rule{0pt}{3.0in}\rule{\linewidth}{0pt}}
   \caption{Example of a short caption, which should be centred.}
\label{fig:short}
\end{figure}

It is possible to use the full width of the page to display images. It is also possible to combine image parts, as shown in Figure~\ref{fig:teaser}.

\begin{figure*}[htpb]
\centering
\begin{tabular}{ccc}
\fbox{\rule{0pt}{4.2cm}\rule{6.45cm}{0pt}}&
\fbox{\rule{0pt}{4.2cm}\rule{3.83cm}{0pt}}&
\fbox{\rule{0pt}{4.2cm}\rule{6.54cm}{0pt}}\\[2pt]
{\small\color{captionblue}(a)}&{\small\color{captionblue}(b)}&{\small\color{captionblue}(c)}
\end{tabular}
\caption{Example of the use of the full width of the page and labelling
sub-figures as (a), (b), and (c) in the caption.}
\label{fig:teaser}
\end{figure*}

Label tables as ``Table~1'', ``Table~2'', etc. Table captions, by contrast, must appear above the table, also centred and set in 8-pt Times New Roman, as illustrated in Table~\ref{tab:example}.

\begin{table}[htpb]
\centering
\caption{Example of a table caption, which should be centred.}
\label{tab:example}
{\renewcommand{\arraystretch}{1.43}%
\begin{tabular}{|l|c|}
\hline
\textbf{Method} & \textbf{Frobnability} \\
\hline
Theirs & Frumpy \\ \hline
Yours  & Frobbly \\ \hline
Ours   & Makes one's heart Frob \\
\hline
\end{tabular}}
\end{table}

Each figure and table must be explicitly cross-referenced within the body of the paper prior to its appearance. Do not include visual elements that are not discussed in the text.



\subsection{Mathematics}

Please number all equations. It is important for readers to be able to refer to each equation.

%-------------------------------------------------------------------------
\section{Bibliographical references}

The entries in the References section should follow an IEEE-like style.

%------------------------------------------------------------------------
\section{Colour}

Colour is encouraged and will be visible to readers of the electronic version. However, authors should ensure that visual elements remain accessible to colour-blind readers and that no important information is lost when printed or converted to grayscale.

%------------------------------------------------------------------------
\section{Lorem ipsum}

Please disregard the content of this section. It is just dummy text to help demonstrate a possible extension for a paper.

Dolor sit amet, consectetuer adipiscing elit. Maecenas porttitor congue massa. Fusce posuere, magna sed pulvinar ultricies, purus lectus malesuada libero, sit amet commodo magna eros quis urna.

Mauris et orci. Aenean nec lorem. In porttitor. Donec laoreet nonummy augue. Suspendisse dui purus, scelerisque at, vulputate vitae, pretium mattis, nunc. Mauris eget neque at sem venenatis eleifend. Ut nonummy.

Suspendisse dapibus lorem pellentesque magna. Integer nulla. Donec blandit feugiat ligula. Donec hendrerit, felis et imperdiet euismod, purus ipsum
pretium metus, in lacinia nulla nisl eget sapien. Donec ut est in lectus
consequat consequat.

Proin nec augue. Quisque aliquam tempor magna. Pellentesque habitant morbi tristique senectus et netus et malesuada fames ac turpis egestas. Nunc ac magna. Maecenas odio dolor, vulputate vel, auctor ac, accumsan id, felis.

%------------------------------------------------------------------------
\section{Also lorem ipsum}

As in the previous section, please disregard the content of this section. It is just dummy text to help demonstrate a possible extension for a paper.

Dolor sit amet, consectetuer adipiscing elit. Maecenas porttitor congue massa. Fusce posuere, magna sed pulvinar ultricies, purus lectus malesuada libero, sit amet commodo magna eros quis urna. Nunc viverra imperdiet enim. Fusce est. Vivamus a tellus. Pellentesque habitant morbi tristique senectus et netus et malesuada fames ac turpis egestas. Proin pharetra nonummy pede.

Mauris et orci. Aenean nec lorem. In porttitor. Donec laoreet nonummy augue. Suspendisse dui purus, scelerisque at, vulputate vitae, pretium mattis, nunc. Mauris eget neque at sem venenatis eleifend. Ut nonummy. Fusce aliquet pede non pede.

Suspendisse dapibus lorem pellentesque magna. Integer nulla. Donec blandit feugiat ligula. Donec hendrerit, felis et imperdiet euismod, purus ipsum pretium metus, in lacinia nulla nisl eget sapien. Donec ut est in lectus consequat consequat. Etiam eget dui. Aliquam erat volutpat. Sed at lorem in nunc porta tristique.

Proin nec augue. Quisque aliquam tempor magna. Pellentesque habitant morbi tristique senectus et netus et malesuada fames ac turpis egestas. Nunc ac magna. Maecenas odio dolor, vulputate vel, auctor ac, accumsan id, felis. Pellentesque cursus sagittis felis. Pellentesque porttitor, velit lacinia egestas auctor, diam eros tempus arcu, nec vulputate augue magna vel risus. Cras non magna vel ante adipiscing rhoncus.

Vivamus a mi. Morbi neque. Aliquam erat volutpat. Integer ultrices lobortis eros. Pellentesque habitant morbi tristique senectus et netus et malesuada fames ac turpis egestas. Proin semper, ante vitae sollicitudin posuere, metus quam iaculis nibh, vitae scelerisque nunc massa eget pede. Sed velit urna, interdum vel, ultricies vel, faucibus at, quam. Donec elit est, consectetuer eget, consequat quis, tempus quis, wisi.


\bibliography{RECPAD2026_template}

\end{document}
